\section{Introduction}\label{sec:intro}

When a flood covers a piece of land that secures a bank loan, three things can happen to the loan. The borrower may stop paying; the collateral may be worth less if the loan is enforced; and the bank may, or may not, recognise either change in its provisions and capital. The evidence on this chain comes almost entirely from the United States and Western Europe, where flood insurance and public aid pay for much of the damage and where default rises only modestly after a flood \citep{gallagher2017household,albertazzi2026physical}. In most emerging markets neither condition holds. Insurance covers a small fraction of flood losses \citep{swissre2022extreme}, and the buffer that stands between a flooded borrower and default is regulatory and fiscal relief: classification freezes, moratoria and subsidised credit. Relief changes not only how large the loss is but when and where it appears in a bank's accounts.

This paper studies the chain in Thailand. The 2011 Chao Phraya floods caused damage and losses of THB~1.43 trillion, around 90\% of it borne by the private sector, while fewer than 1\% of private dwellings were insured against natural disasters \citep{worldbank2012thai}. Floods have recurred since, in the north in 2024 and in the north-central and southern provinces in 2025. The Bank of Thailand's response in 2011 was a circular that let lenders keep the pre-flood classification of performing borrowers hit by the floods and defer collateral revaluation \citep{bot2011circular1564}; since 2024 a standing circular, framed as a request for cooperation, has asked institutions to apply the same template after any declared disaster or alert \citep{bot2024circular601}. Three features make Thailand a useful laboratory. The insurance channel that dominates the advanced-economy evidence is largely switched off. Relief is rule-based, with eligibility cut-offs that can be used for identification. And the state publishes, as open data, the official value of every titled land parcel, reset in four-year rounds on a calendar that does not depend on floods.

The paper has three parts, which correspond to three questions. \emph{Question~1} asks whether the state's official land values respond to realised floods, and, with an outside estimate of the flood discount, how much of it they capitalise. \emph{Question~2} asks whether flood exposure of collateral raises borrower distress and loss given default, and whether under relief the effect surfaces first in restructuring and in IFRS~9 stage~2 rather than in 90-day default. \emph{Question~3} asks what flood severity would breach a bank's capital threshold once observed exposure, valuation behaviour and relief are taken into account: a reverse stress test over a physical hazard variable.

The three questions are at different stages, and the paper says so at every step. It reports one descriptive result, one pre-specified design whose estimate awaits the next valuation round, and tested code for the credit and capital steps. For Question~1 the descriptive result is a cross-section of the 2023--26 valuation round against satellite-measured 2011 inundation across \pilotParcelsM~million parcels, which I argue does not identify capitalisation; the design is a difference-in-differences comparison of the next valuation round with the current one on sheets flooded in 2024 and 2025, for which the radar flood maps, validation, control selection, balance, power and Monte Carlo evidence are complete and the treatment and control lists are hashed. Registration of the plan is pending; the estimate cannot be produced until the Treasury Department publishes the next round, which as of 4 November 2026 is likely to be delayed by a year to 1~January 2028 (Section~\ref{sec:status}). For Question~2 the loan-level data are restricted, and the paper reports what public aggregates show about credit after 2011, the estimator and the power calculations, with the recovery of planted effects on a synthetic panel left to the appendix. For Question~3 the stress-test machinery is complete and is exercised on the synthetic panel and on an illustrative calibration; its empirical inputs await Questions~1 and~2.

Four things can be stated now. First, in the Chao Phraya basin, official land values on urban cadastral sheets that were almost wholly flooded in 2011 are \binTopUrbanPct\% below values on sheets in the same province and 1:50,000 map block that stayed dry, and the gradient is monotone across flood-share bins. Holding land cover fixed removes most of it: flooded urban sheets are \lcUrbanBuiltTop\% built-up against \lcUrbanBuiltRef\% for dry ones, and the residual urban gap is about 8\% and imprecise, with land cover measured ten years after the flood. On rural sheets the raw gradient is weaker and turns positive with controls. Because flood-prone land differs from dry land in use, elevation and amenity, these are not estimates of the flood discount; they show that the official values carry a large amount of flood-correlated variation and that a cross-section cannot tell hazard from land use. Second, Sentinel-1 radar flood maps of the 2024 northern floods classify \radarTreated{} cadastral sheets in four provinces as treated; against Sentinel-2 optical imagery, \valControlContam\% of the chosen control sheets show flooding and \valTreatedPrec\% of treated sheets do, so if misclassification is unrelated to what the committees know about hazard it attenuates the estimate. Most treated sheets are rice land. Third, with \radarTreatedClusters{} treated 1:50,000 clusters the design has a minimum detectable effect of \regMdePooled{} log points at 80\% power under the registered weights (standard deviation of value changes 0.10, intra-cluster correlation 0.2), about the size of a fully recognised 4.6\% hazard discount \citep{beltran2018flood}; its power against full recognition of that discount is \powFullMid, against half recognition \powHalfMid, and against realised-damage discounts of 10\% or more above \powDamageTen. In \mcxReps{} Monte Carlo replications on synthetic next rounds the wild cluster bootstrap rejects a true null \mcxNullWcb{} of the time, but the placebo-footprint test, which the draft plan made primary, rejects it \mcxNullRi{} of the time, because shifted footprints land on better-overlapped sheets than the real one and are estimated too precisely; restricting where placebo footprints may land does not repair the test, studentising the comparison does (size \prxNullRegT), and the plan is amended before registration to make the bootstrap primary and the placebo check studentised (Section~\ref{sec:montecarlo}). A larger 2025 sample has a minimum detectable effect of \evMdeW. Fourth, in an illustrative stress test, whether valuations price hazard moves the breaking point: at a loan-to-value ratio of 0.8 and a 4.6\% discount, moving from valuations that ignore hazard to valuations that price it fully raises the flood severity that breaches capital by \lamSStarEightChange\%, and by \lamSStarOneChange\% at a ratio of 1.0; at lower ratios the effect is larger but depends on the assumed damage share and haircut.

The closest work is as follows. \citet{guo2026rising} show that after Hurricane Sandy the assessed market values of inundated New York City homes fell by about 17\%, mostly through tax relief, while those of undamaged homes in the high-risk zone did not change even though their sale prices fell, and that statutory assessment caps blunt the pass-through; Question~1 is the emerging-market counterpart of that test, nationwide in coverage and with a satellite-measured treatment, and its design is the paper's main contribution. \citet{garbarino2021high} show that after the 2013--14 English floods sale prices fell but lenders' valuations did not, an upward bias of about 3\%; that finding is the $\lambda_B$ mechanism this paper assumes rather than measures, and \citet{vangerwen2025flooded} report the opposite behaviour at one Belgian bank after the 2021 floods, which revalued exposed collateral downward. \citet{albertazzi2026physical} combine satellite flood perimeters with euro-area credit-register data and find that default on existing loans rises by about 0.6--0.7 percentage points on a base of about 1.4\%, a near doubling, with a credit surge that reverses and a rise in loan-to-value ratios on real-estate-backed loans to micro firms; their setting is an advanced economy with partial insurance. \citet{perezmontes2025impact} show that after the 2024 Valencia floods stage~2 shares rose by about half a percentage point and that the most exposed banks moved more loans into stage~2 and fewer into stage~3 while moratoria were in force. \citet{blickle2025bad} compute the disaster severity that would fail a single-county US bank, in effect a reverse stress test on a hazard variable, and find banks resilient; \citet{caloia2021flood} and \citet{caloia2026floods} run Dutch banks' geocoded collateral through fixed flood scenarios and find non-linear capital depletion. Section~\ref{sec:stress} builds on these by solving the inversion for a scaling of the flood footprint, with relief and valuation behaviour as explicit inputs. In Thailand, \citet{chitchumnong2025physical} study one bank's mortgage book around 2011 and find default below model-predicted levels after the flood, which they attribute to relief, with no geolocation, collateral values or staging; \citet{sawada2018land} find that industrial land prices fell by 6--8\% after 2011, while residential and commercial land showed no significant change; and \citet{chongwilaikasaem2022flood} find a negative effect of flood risk on official valuations and rents in Hat Yai. No study yet traces observed inundation of individual collateral through official and bank valuations to default, restructuring, staging and capital where relief rather than insurance is the buffer. This paper sets out how to do so and takes the first, public-data step; it does not yet trace the chain.

The paper also draws on three wider literatures. On the pricing of hazard in property markets, exposure to sea-level rise carries a discount that grows with buyer sophistication and beliefs \citep{bernstein2019disaster,baldauf2020does}, sales volumes fall before prices do \citep{keys2020neglected}, lenders ration uninsurable flood risk through loan-to-value limits \citep{sastry2026who}, and official assessments are less sensitive to neighbourhood attributes than prices are \citep{avenancioleon2022assessment}; flood-exposed urban land, moreover, keeps its economic activity after floods \citep{kocornikmina2020flooded}. On debt relief, waivers and moratoria reallocate credit, change later repayment and can conceal weak loans \citep{gine2018economic,mukherjee2018borrowers,dinerstein2024debt,kim2024intermediation,chari2021unholy}, which is why Question~2 treats relief as a censoring process rather than a transfer. On flood stress tests, capital effects are non-linear in damage and sensitive to defences and to collateral modelling \citep{lepore2024beyond,dhima2026integration}, and in Thailand the systemic-risk evidence so far concerns transition risk, with physical risk flagged as the open question \citep{luangaram2024climate}; household evidence from 2011 finds income losses with small effects on debt \citep{noy2021floods}.

One framing point governs the interpretation throughout. The official values studied in Question~1 are the state's base for transfer fees and taxes, set by provincial committees by market comparison; they are not banks' collateral values, and no public Bank of Thailand rule links the two \citep{th2020appraisalmr,bot2019collateral}. Question~1 therefore measures the first-round response of the official valuer, and its relation to the bank parameter $\lambda_B$ of Section~\ref{sec:framework} is a stated assumption that loan-level appraisal data would test.

Section~\ref{sec:background} sets out the institutions. Section~\ref{sec:framework} gives the framework and its testable predictions. Section~\ref{sec:data} describes the data. Section~\ref{sec:cross} reports the 2011 cross-section. Section~\ref{sec:design} reports the pre-specified design for Question~1 and its validation, power, timing and status. Section~\ref{sec:credit} covers Question~2. Section~\ref{sec:stress} covers Question~3. Section~\ref{sec:discussion} discusses what can and cannot yet be concluded, and Section~\ref{sec:conclusion} concludes. Appendices~\ref{app:tables} to~\ref{app:timing} hold the supporting tables, the code tests on synthetic data and the timing cases.
